\section{Transport of the radial covariance deficit}\label{cap:multiscale}
We now propagate a covariance estimate from one cap to the next. The
spherical saddle supplies a radial deficit, while a fresh Haar refinement
reduces the transverse error. Section~\ref{cap:initialization} constructs an
initial potential and replaces it, after finitely many levels, by the true
partition function. Work on successive levels \(p,c=\gamma p\) with
\(\gamma\) a very small fixed ratio. For \(s>0\), define
\begin{equation}
 \mathcal A_p(s)=\{h\in H_p:|\alpha_p(h)|\le s p^2\};
 \qquad\text{the induction uses }\mathcal A_p(\sigma)
 \label{cap:annulus}
\end{equation}
with small fixed \(\sigma>0\). We will start with auxiliary smooth log
potentials close in \emph{gradient norm by a bounded amount} to the true ones
there. This auxiliary phase lasts only a fixed number of levels, which can be
arbitrarily large but is independent of \(n\). Thereafter the potentials are
the true ones. In either case write
\[
 K_p=\nabla_{H_p}^2(\phi_p|_{H_p}),\qquad
 R_p^H=(P_pR_pP_p)|_{H_p}=(p^2+\alpha_p(h))^{-1}I_{H_p}.
\]
For the true potential,
\(K_p=\operatorname{Cov}_{\mu_{p,h}}(P_pX)\) as an operator on \(H_p\).
Thus every cap curvature assertion below concerns \(H_p\). In
the auxiliary case we require smoothness only on the used annulus, and the
starting auxiliary function can itself be a Gaussian convolution. Potentials
at subsequent levels in the auxiliary phase are computed by heat convolution,
using a semiconcave extension at each step, defined below. Additive
calibrations of these potentials on \eqref{cap:annulus} are harmless. The
invariants are as follows; all uniformities on an annulus are pointwise
everywhere. The saddle equation gives
\[
 \|h\|=(p^2+\alpha)\sqrt{n-\Tr R_{p,\alpha}}
       =(\sqrt b+O(\sigma)+o(1))p^2M_p
       \quad\text{when }|\alpha|\le\sigma p^2 .
\]
The annulus is therefore separated from zero. In the frame radial and
transverse to \(e=h/\|h\|\), put
\(D=p^2(K_p-R_p^H)\in\operatorname{Sym}(H_p)\). Then
\begin{equation}
 \|D_{\perp\perp}\|\le\delta,\quad
 \|D_{\perp e}\|\le K_0\sqrt\delta,\quad
       -K_0 I_{H_p}\preceq D,\quad D_{ee}\le-d_0,\qquad d_0=0.30 .
 \label{cap:invariant}
\end{equation}
Here \(K_0\) is a sufficiently large absolute constant. We can take \(\delta\)
a sufficiently high fixed power of \(\gamma\), and \(\sigma\) likewise.
The negative radial sign also controls the full positive part. If
\(z=D_{\perp e}\) and \(v=v_\perp+v_e e\), then
\[
 v^\top Dv\le\delta\|v_\perp\|^2
 +2|v_e|\|z\|\|v_\perp\|-d_0v_e^2
 \le\left(\delta+\frac{\|z\|^2}{d_0}\right)\|v_\perp\|^2.
\]
Hence \(D\preceq C\delta I_{H_p}\) and
\(\|D\|\le\max(K_0,C\delta)\).
Throughout multiscale estimates constants polynomially bad in \(\gamma^{-1}\)
from inverses with fixed spectral formulas are allowed (the powers do not
depend on \(\delta,\sigma\)); the small powers used for tolerances can be
chosen with sufficiently large fixed exponents. Macroscopic value and gradient
errors relative to the spherical predictions can be taken as tiny as
subsequently required. This is by \eqref{cap:static-values}; in the auxiliary
phase we use the gradient comparison just mentioned with the true model, which
integrates to \(o(k_p)\) value oscillation error on \eqref{cap:annulus} since
the scales in that phase are fixed constants. The starting potential, its
auxiliary randomness if any, and the gradient error are intrinsic given the
current cap data and independent of the unused Haar refinement.

The extension used in an auxiliary step is one function selected from
the full parent annulus using only these parent data. It agrees with the
parent potential on an inner annulus with fixed slack, has at most
quadratic growth, and has a global upper Hessian bound
$(1+\gamma^4/100)I/p^2$ after the curvature tolerances are made sufficiently
small. This is the semiconcavity bound that controls its Gaussian
convolution. The supporting-quadratic construction below verifies these
properties without using the child eigenframe. For a true parent the
extension is used only to estimate the transition; the output remains
the true child potential.

\begin{proposition}[One cap step]\label{cap:curvature-induction}
Fix a bounded child field range for the one-off evaluation conclusion below.
There are a sufficiently large absolute $K_0$ and a sufficiently small
ratio $\gamma>0$ with the following property. Choose $\sigma$ as a
sufficiently high fixed power of $\gamma$, and then choose $\delta$ as a
still higher fixed power, allowing these tolerances to depend on the
prescribed field range. Let a parent potential at scale $p$ be intrinsic
to the capped data, satisfy the value and gradient comparisons above, and
satisfy \eqref{cap:invariant} at all fields in \eqref{cap:annulus}.
Use the semiconcave extension just described when forming an auxiliary
child potential. For a true parent, retain the true child potential.
Condition on the parent data and used auxiliary randomness
\(\mathscr F_p\) described below, on which these hypotheses and the
\(\mathscr F_p\)-measurable parent-level conclusions of the static
comparison hold, including the projected posterior tail and moment bounds
for this capped parent. The separate \(o(1)\) failure probability used to
obtain these parent conclusions remains outside the conditional estimate
below. Then,
uniformly over such parent data, the probability over the unused Haar
refinement \(O\) that \eqref{cap:invariant} fails at \(c=\gamma p\) anywhere
on the child annulus is at most
\[
 C e^{-a k_p}+C e^{-a' k_c}
\]
for fixed positive \(a,a'\). The same inequalities hold with strict slack
on the complementary event. These failure probabilities are summable over
the geometric cap list.

At intermediate ratios $c/p\in[\gamma,1)$, if the frozen parameter stays
inside the agreement domain with fixed slack, the transverse error and
mixed block tend to zero with the parent errors, the total scaled Hessian
is bounded, and the radial deficit is at least $0.26$. The same deterministic
evaluation bounds hold for $c/p\in[\gamma^2,\gamma]$ under this
interior-parameter condition. For a true parent and the prescribed bounded
child field range there is also a sufficiently small fixed evaluation-ratio
interval on which the upper child curvature is at most $C/c^2$, with $C$ depending only on
that field range. This last evaluation uses supports centered at the frozen
parameter, constructed at the end of the proof.
All statements are uniform down to $np^3\asymp\sqrt{\log n}$.
\end{proposition}
The proof occupies the remainder of this section. Its probabilistic step is
used only at ratio $\gamma$; the intermediate-scale conclusions are
consequences of the deterministic covariance bounds.

\subsection{Semiconcave extensions and the observation bridge}
We make the measurability in a refinement explicit. Conditional on the
eigenvalues and the capped matrix $J_p$, let $\mathscr F_p$ include all
auxiliary randomness already used to construct the parent potential. The
construction is required to be measurable with respect to $\mathscr F_p$;
it does not reveal the orientation of the unequal eigenvalues within the
flat eigenspace $P_p$. Choose a fixed orthonormal identification of this space
with $\mathbb R^{d_p}$, where $d_p=\Tr P_p$. The remaining refinement is an
orthogonal matrix $O$ which is Haar and independent of $\mathscr F_p$.
All spectral projectors and increments below are obtained by conjugating
fixed diagonal matrices by $O$. This independence is the reason for the
auxiliary refinement used in Section~\ref{cap:initialization}.

After a step, the increment $\Delta=B_p-B_c$ is a function of $J_p$ and $J_c$
and is scalar on $P_c$. Thus the output auxiliary potential depends only on
the new capped data and previously introduced independent seeds; it does not
reveal the internal eigenframe of $P_c$. Conditional on these data the next
finer rotation is Haar.

For \(c=xp<p\), put
\[
 \Delta=B_p-B_c,\qquad E=\operatorname{ran}\Delta\subset H_p,\qquad
 U=H_c\subset E.
\]
Every inverse of \(\Delta\) is taken on \(E\), after omitting the zero
increment coordinates. An inverse is taken before its block is compressed
to \(U\). In particular \([M^{-1}]_{UU}\) generally differs from
\([M_{UU}]^{-1}\). The child increment is the scalar
\(\Delta_U=(1-x^2)p^2I_U\).

For a child field \(h_c\), the exact convolution and its transition
probability are
\begin{align}
 \phi_c(h_c)&=\log\int \exp\{\phi_p(h_c+z)\}\,\mathsf N(0,\Delta)(\dd z),
 \label{cap:heat-convolution}\\
 \mathsf T_{c,p}^{h_c}(\dd h_p)
 &=\frac{\exp\{\phi_p(h_p)\}\,
       \mathsf N(h_c,\Delta)(\dd h_p)}{\exp\{\phi_c(h_c)\}}.
 \label{cap:actual-transition}
\end{align}
For an auxiliary step, the function convolved in these equations is the
selected extension \(\bar\phi_p\) defined below; its logarithmic convolution
defines the child potential. For the true partition function the equations
follow by integrating the Gaussian linear tilt for each spin. Freeze the child
saddle parameter \(\alpha\) and put \(R_u=R_{u,\alpha}\) for \(u=c,p\).
Replacing the convolved potential by \(q_{p,\alpha}^{\rm fr}\) from
\eqref{cap:frozen-potential} gives the canonical probability
\(\mathsf G_{c,p}^{h_c}\). Under this law,
\begin{equation}\label{cap:canonical-bridge}
 h_p=h_c+\Delta a_c+\zeta,\qquad
 \zeta\sim\mathsf N(0,\Delta+\Delta R_c\Delta),\qquad
 R_p h_p\sim\mathsf N(a_c,R_c-R_p).
\end{equation}
These identities follow from \(R_c^{-1}=R_p^{-1}-\Delta\), as in
\eqref{cap:static-canonical-bridge}. The covariance \(R_c-R_p\) is supported
on \(E\). The reference spin Gaussian is sampled without norm conditioning.

The parent Hessian estimate holds on an annulus, whereas the convolution
integrates over the full increment space. We extend the potential to make the
convolution globally controlled, while preserving it on the fields the bridge
visits with high probability. The extension must be determined by the parent
data alone: using the child refinement in its construction would consume the
randomness needed for transverse compression.

Here is one parent-measurable extension with the needed buffers. The radius
corresponding to a saddle parameter is
\[
 \ell_p(\alpha)=(p^2+\alpha)\sqrt{n-\Tr R_{p,\alpha}},\qquad
 \ell_p'(\alpha)=\sqrt{n-\Tr R_{p,\alpha}}+
 \frac{(p^2+\alpha)\Tr R_{p,\alpha}^2}
      {2\sqrt{n-\Tr R_{p,\alpha}}}\asymp M_p .
\]
Thus \(\mathcal A_p(s)\) is the radial annulus between
\(\ell_p(-sp^2)\) and \(\ell_p(sp^2)\). Put
\[
 S_0=\mathcal A_p(\sigma/2),\qquad
 S_1=\mathcal A_p(3\sigma/4),\qquad
 S_2=\mathcal A_p(\sigma).
\]
Half the distance from \(S_0\) to \(H_p\setminus S_1\), and from \(S_1\)
to \(H_p\setminus S_2\), gives buffers \(\rho_{01},\rho_{12}\asymp
\sigma p^2M_p\). In particular \(S_0+B(0,\rho_{01})\subset S_1\),
and a segment of length at most \(\rho_{12}\) from a point of \(S_1\)
lies in \(S_2\).

Let \(f_p=\phi_p|_{H_p}\) be the parent before extension and set
\[
 L_p=(1+\zeta)p^{-2}I_{H_p},\qquad \zeta=\gamma^4/100.
\]
We claim that every pair \(z,w\in S_1\) satisfies
\begin{equation}\label{cap:parent-upper-support}
 f_p(w)\le f_p(z)+\nabla_{H_p}f_p(z)^\top(w-z)
                  +\tfrac12(w-z)^\top L_p(w-z).
\end{equation}
For \(\|w-z\|\le\rho_{12}\), the segment stays in \(S_2\). There
\eqref{cap:invariant} and its positive-part bound give
\[
 K_p\preceq p^{-2}\bigl((1-\sigma)^{-1}+C\delta\bigr)I_{H_p}
 \prec L_p
\]
after choosing \(\sigma,\delta\) below the excess \(\zeta\). Integrating the
Hessian on the segment proves \eqref{cap:parent-upper-support} in this case.

For \(\|w-z\|\ge\rho_{12}\), put \(d=w-z\) and freeze the own saddle at
\(z\). The value comparison, the gradient comparison, and
\eqref{cap:frozen-mgf} give
\[
 f_p(w)-f_p(z)-\nabla_{H_p}f_p(z)^\top d
 \le \frac{\|d\|^2}{2(p^2+\alpha_p(z))}
       +2\varepsilon_v k_p+C_{\rm dens}
       +\varepsilon_g M_p\|d\|.
\]
In an auxiliary step, \(\varepsilon_v,\varepsilon_g\) also include the
calibrated auxiliary errors. With
\(\beta_{\rm buf}=\rho_{12}/(p^2M_p)\), the last three terms are at most
\[
 \frac{\|d\|^2}{p^2}
 \left(\frac{2\varepsilon_v+C_{\rm dens}/k_p}{\beta_{\rm buf}^2}
       +\frac{\varepsilon_g}{\beta_{\rm buf}}\right).
\]
Choose these errors after the buffers so that this quantity fits below the
remaining excess of \(L_p\). This proves
\eqref{cap:parent-upper-support} for every support pair.

Now define the single function
\begin{equation}\label{cap:parent-extension}
 \bar\phi_p(y)=\inf_{z\in S_1}
 \left\{f_p(z)+\nabla_{H_p}f_p(z)^\top(y-z)
                   +\tfrac12(y-z)^\top L_p(y-z)\right\}.
\end{equation}
At \(y\in S_1\), \eqref{cap:parent-upper-support} bounds every term below
by \(f_p(y)\), and the term \(z=y\) equals it. Thus \(\bar\phi_p=f_p\)
on \(S_1\), in particular on a neighborhood of \(S_0\). Subtracting
\(\tfrac12y^\top L_py\) leaves an infimum of affine functions, hence a
concave function. The support set is bounded and its values and slopes are
bounded, so this infimum is finite and has at most quadratic growth.
Consequently \(\nabla^2\bar\phi_p\preceq L_p\) in the distributional sense.
The construction uses only the parent annulus and parent function, so it is
\(\mathscr F_p\)-measurable and consumes no part of \(O\).

On \(E\) the convolution has the strict margin
\[
 \|\Delta_E^{1/2}L_p\Delta_E^{1/2}\|
 \le(1+\zeta)(1-x^2)
 \le(1+\zeta)(1-\gamma^2)<1 .
\]
The extended terminal density is therefore strongly log-concave in the
whitened increment coordinates. In the true phase the extension is used only
to estimate the transition; the output is the true child. If smoothing is
needed, mollify the extension and pass to the limit at fixed \(n\).
Inside the buffered agreement region the mollifications converge to the
original smooth potential with its derivatives.

We next record precisely how canonical estimates pass to a terminal
transition. Let \(V_p\) be the potential actually convolved, either
\(\bar\phi_p\) in an auxiliary step or the true \(\phi_p\), and set
\[
 F(y)=V_p(y)-q_{p,\alpha}^{\rm fr}(y).
\]
For fixed parent data and child field, the exact likelihood ratio is
\begin{equation}\label{cap:bridge-density-ratio}
 \frac{d\mathsf T}{d\mathsf G}(y)
 =\frac{e^{F(y)}}{\E_{\mathsf G}e^F},
 \qquad \mathsf G=\mathsf G_{c,p}^{h_c}.
\end{equation}
Let \(\widehat\alpha=\alpha_p(y)\) and take a shell
\(\mathcal S_\tau=\{|\widehat\alpha-\alpha|\le\tau p^2\}\) lying inside
the buffered agreement region. Under \(\mathsf G\), its complement has
probability at most \(e^{-a_\tau k_p}\) for each fixed \(\tau>0\), by
\eqref{cap:canonical-bridge} and the saddle derivative estimate. On
\(\mathcal S_\tau\), the static comparison or its calibrated auxiliary
version, together with \eqref{cap:off-saddle-density}, gives, for a separately
chosen normalization shell \(\mathcal S_{\tau_0}\) with \(0<\tau_0<\tau\),
\begin{equation}\label{cap:shell-likelihood}
 \sup_{\mathcal S_\tau}\frac{d\mathsf T}{d\mathsf G}
 \le \mathsf G(\mathcal S_{\tau_0})^{-1}
       \exp\{[2\varepsilon_v+C(\tau^2+\tau_0^2)]k_p+O_x(1)\}.
\end{equation}
The denominator is bounded below by integrating \(e^F\) on the narrower
shell. Choosing \(\tau\), then \(\tau_0\), and then the comparison errors makes the exponent
at most \(\eta k_p\) for any prescribed fixed \(\eta>0\).

The shell complement requires its own estimate. Write
\(d(y)=|\|y\|-\ell_p(\alpha)|/(p^2M_p)\). The canonical bridge has
\[
 \E_{\mathsf G}\|Y\|^2=\ell_p(\alpha)^2,
 \qquad \|\Cov_{\mathsf G}(Y)\|\le C\gamma^{-2}p^2.
\]
The first identity follows from
\(\E\|R_pY\|^2=\|a_c\|^2+\Tr(R_c-R_p)=n-\Tr R_p\),
since \(R_p\) is scalar on \(H_p\). Gaussian concentration for the norm,
together with
\(0\le\sqrt{\E\|Y\|^2}-\E\|Y\|\le
\sqrt{\|\Cov(Y)\|}\), therefore gives
\begin{equation}\label{cap:bridge-distance-tail}
 \Pr_{\mathsf G}\{d(Y)>v\}\le 2e^{-c\gamma^2k_pv^2},
 \qquad v\ge \frac{C}{\gamma\sqrt{k_p}}.
\end{equation}
In particular this estimate applies at every fixed shell threshold once
\(n\) is sufficiently large. For the extended potential, a support from
the frozen typical sphere, with the calibrated value there,
bounds \(F(y)\) above by a fixed reference level plus
\[
 \eta_{\rm err}k_p+\eta_{\rm small}k_p d(y)^2.
\]
Leaving \(\mathcal S_\tau\) requires
\(d(Y)\ge c_0\tau\) on the comparable-gap range. Choose
\(\eta_{\rm small}<c\gamma^2/2\), then choose \(\tau_0\) and the comparison
errors so that the normalization loss
\(2\varepsilon_v+C\tau_0^2\), together with \(\eta_{\rm err}\), is less than
\((c\gamma^2-\eta_{\rm small})c_0^2\tau^2/2\).
The narrower shell has canonical probability tending to one for every fixed
\(\tau_0>0\). Integrating the tail against the upper support with this
normalization budget
gives, for every fixed integer \(j\),
\[
 \E_{\mathsf T}\bigl[(1+d(Y))^j\mathbf1_{\{Y\notin\mathcal S_\tau\}}\bigr]
 \le C_{j,\tau}e^{-a'_\tau k_p}.
\]
For a true transition, on each bounded field range the static value
comparison and the frozen numerator norm-density upper bound control
\(F\) by its calibrated level plus \(\varepsilon_vk_p+C\). The same
normalization budget and \eqref{cap:bridge-distance-tail} then control the
bounded part of the shell complement. Beyond that range, the projected \(v^6\) tail at
parent zero field is multiplied by at most \(e^{Ck_p(1+v^2)}\) from the
added quadratic, so \(e^{-ck_pv^6+Ck_p(1+v^2)}\) bounds the large-length
contribution and all fixed polynomial moments. Gaussian observation noise
has the corresponding quadratic tail. These estimates make the error from
replacing a true transition by its extension exponentially small in scaled
covariance units. Indeed the two unnormalized terminal densities agree on
\(S_1\). If each normalized law assigns at most \(\varepsilon\) to its
complement, their normalization ratio lies between \(1-\varepsilon\) and
\((1-\varepsilon)^{-1}\); the complementary moment bounds then control the
difference of their means and covariances.

The same statements hold for \(x\in[\gamma,1)\) whenever the frozen
parameter lies inside parent agreement with fixed slack; one use is
\(|\alpha|\le2\sigma(\gamma p)^2\). As \(x\uparrow1\), the canonical
increment covariance shrinks proportionally to \(1-x\), so the probability
of leaving agreement is at most \(C e^{-a k_p/(1-x)}\). The support and
projection estimates give the same bound for the relevant true-transition
errors. This exponential factor pays every fixed power of \(1/(1-x)\).
For derivative integrals on a true transition we use the projected trace
bound
\[
 \Tr_{H_p}K_p(y)
 \le \E_{\mu_{p,y}}\|P_pX\|^2.
\]
The projected moment tails therefore leave only polynomial factors in
\(k_p\), rather than an ambient rank factor.

\subsection{Upper and lower covariance transport}
We use the Brascamp--Lieb variance inequality~\cite{BrascampLieb1976};
its relation to log-concavity and logarithmic Sobolev inequalities is
discussed in~\cite{BobkovLedoux2000}. The paired upper and lower covariance
bounds also have the form used in~\cite[Lemma~7 and Appendix~A]{ChewiPooladian2023Entropic}.
The output Hessian is obtained from the covariance of the terminal field.
We bound this covariance from both sides. Brascamp--Lieb keeps the inverse
curvature correction and will preserve the useful negative radial term.
The Fisher test is linear in the parent correction and supplies the lower
bound needed for transverse and mixed blocks. We perform the lower test with
cutoffs inside agreement, so uncontrolled curvature of the extension is not
integrated.

Differentiating the logarithmic convolution in a direction of \(U\) gives
the exact output Hessian for the potential \(V_p\) and its corresponding
transition \(\mathsf T\):
\begin{equation}\label{cap:child-hessian}
 K_c=\Delta_U^{-1}
        [\operatorname{Cov}_{\mathsf T}(Y)]_{UU}
        \Delta_U^{-1}-\Delta_U^{-1}.
\end{equation}
The covariance is first computed on \(E\). For Brascamp--Lieb, temporarily
take a smooth mollification \(V_p^{\rm ext}\) of the extension and denote
its transition by \(\mathsf T^{\rm ext}\). Put
\[
 K_p^{\rm ext}(Y)=\nabla^2 V_p^{\rm ext}(Y),\qquad
 M_{\rm ext}(Y)=\Delta_E^{-1}-[K_p^{\rm ext}(Y)]_{EE},\qquad
 C_0=(\Delta_E^{-1}-R_p|_E)^{-1}
       =\Delta+\Delta R_c\Delta .
\]
The extension margin makes \(M_{\rm ext}\) positive. Brascamp--Lieb on \(E\) gives
\(\operatorname{Cov}_{\mathsf T^{\rm ext}}(Y)
\preceq\E_{\mathsf T^{\rm ext}}M_{\rm ext}(Y)^{-1}\).
Combining this with \eqref{cap:child-hessian} preserves the inverse before
compression to \(U\).

For the lower estimate under this smooth extension law, let
\(\rho_{\rm ext}\) be its density and
\(\mathfrak s=-\nabla_E\log\rho_{\rm ext}\) be the
score. Integration by parts gives
\(\E[(Y-\E Y)\mathfrak s^\top]=I_E\) and
\(\E[\mathfrak s\mathfrak s^\top]=\E M_{\rm ext}\). The nonnegativity of the covariance
of \(Y-T\mathfrak s\) therefore gives, for every constant matrix \(T\) on \(E\),
\begin{equation}\label{cap:fisher-variational}
 \operatorname{Cov}_{\mathsf T^{\rm ext}}(Y)
 \succeq T+T^\top-T\,\E_{\mathsf T^{\rm ext}}[M_{\rm ext}(Y)]\,T^\top.
\end{equation}
With \(T=C_0\), the right-hand side is
\(C_0+C_0\E_{\mathsf T^{\rm ext}}[K_p^{\rm ext}-R_p]_{EE}C_0\), which is linear in the
parent correction.

We use a localized version when only the agreement region has controlled
classical curvature. This identity applies separately to either smooth
law with density \(\rho\), using its own score \(\mathfrak s\), precision
\(M=-\nabla_E^2\log\rho\), and expectations. Let \(\chi\) be a smooth cutoff supported there and,
for constant \(t,w\in E\), put
\(b_\chi=\chi\,t^\top\mathfrak s-t^\top\nabla\chi\).
Integration by parts yields
\[
 \E[(w^\top(Y-\E Y))b_\chi]=\E\chi\,w^\top t,\qquad
 \E b_\chi^2=\E\chi^2 t^\top M t+\E(t^\top\nabla\chi)^2.
\]
The inequality \(\E(w^\top(Y-\E Y)-b_\chi)^2\ge0\) becomes
\begin{equation}\label{cap:localized-fisher}
 \operatorname{Var}(w^\top Y)
 \ge2\E\chi\,w^\top t-\E\chi^2t^\top Mt
                       -\E(t^\top\nabla\chi)^2 .
\end{equation}
Choose \(\chi=1\) on a narrower shell and supported inside agreement.
For an output test \(v\in U\), take \(w=\Delta_U^{-1}v\) and
\(t=C_0w=Nv\), where
\(N=(p^2+\alpha)/(c^2+\alpha)\). Only the scalar inverse \(\Delta_U^{-1}\)
appears in this test. The shell estimates above, including their polynomial
moment bounds and the improved \(e^{-a k_p/(1-x)}\) rate, make the cutoff
errors exponentially small after any fixed powers of \(k_p\) and
\((1-x)^{-1}\). Thus \eqref{cap:localized-fisher} applies to the true
transition as well as a smoothed extension without integrating uncontrolled
curvature outside agreement.

On the shell subtract the frozen cap resolvent, writing
\[
 D^\circ(y)=p^2\bigl(K_p(y)-(p^2+\alpha)^{-1}I_{H_p}\bigr).
\]
If \(\widehat\alpha=\alpha_p(y)\), then exactly
\[
 D^\circ(y)-D(y)
 =\frac{p^2(\alpha-\widehat\alpha)}
        {(p^2+\widehat\alpha)(p^2+\alpha)}I_{H_p}.
\]
This is \(O(\tau)\) on \(\mathcal S_\tau\). In the following resolvents
\(D^\circ\) is restricted to \(E\); the restriction is suppressed in the
notation. Put
\[
 S=p^{-2}C_0=p^{-2}\Delta(I-R_p\Delta)^{-1}.
\]
On the agreement shell, where the limiting extension Hessian is the parent
Hessian, put \(M=C_0^{-1}-p^{-2}D^\circ\). The exact identities there are
\[
 M^{-1}=C_0+p^{-2}C_0D^\circ(I-SD^\circ)^{-1}C_0,\qquad
 \Delta_U^{-1}(C_0)_{UU}\Delta_U^{-1}-\Delta_U^{-1}=R_c|_U.
\]
Together with \(C_0\Delta_U^{-1}=N I_U\), these identities give the
extension bound. Pass to the mollification limit and, for a true output,
use the preceding transition comparison. The cutoff and change-of-law
errors are included in \(\eta_{\rm err}\). The resulting target bound is
\begin{equation}\label{cap:bl-recursion}
 c^2(K_c-R_c|_U)
 \preceq x^2N^2\,\E_{\mathsf T}
       [D^\circ(I-SD^\circ)^{-1}]_{UU}
       +\eta_{\rm err}I_U .
\end{equation}
The constant Fisher test gives the
corresponding lower bound \(x^2N^2\E_{\mathsf T}[D^\circ]_{UU}\), up to the
same type of errors. Testing only \(v_\perp\), now transverse to the child
field, gives that transverse lower block with zero mixed block and radial
block \(-c^2(\Delta_U^{-1}+R_c|_U)\). The global upper curvature controls
the rare-event remainder in \eqref{cap:bl-recursion}; all expansions of
\(D^\circ\) use the good shell.

The matrix map in the upper bound preserves Loewner order on its positive
domain. On \(E\),
\begin{equation}\label{cap:resolvent-order}
 D(I-SD)^{-1}
 =S^{-1/2}\bigl[(I-S^{1/2}DS^{1/2})^{-1}-I\bigr]S^{-1/2}.
\end{equation}
It is increasing whenever \(I-S^{1/2}DS^{1/2}\succ0\). The extension
margin supplies this condition for the actual correction. It will also hold
for the upper comparison below: its positive part is \(O(\delta)+
\eta_{\rm err}\), whereas \(\|S\|\) is bounded by a fixed power of
\(\gamma^{-1}\). The stated tolerance order leaves a strict margin.

\subsection{Regeneration of the radial deficit}
Set \(A=D_{\perp\perp}\), embedded with zero radial block, and
\(z=D_{\perp e}\). Completing the square in the mixed block, using a small
fixed part of the radial deficit, gives
\begin{equation}
 D^\circ\preceq A+Czz^\top-d_*ee^\top+\eta_{\rm err}I,
\end{equation}
where \(d_*<d_0\) can be chosen as close to \(d_0\) as desired and
\(C\) depends on \(d_0-d_*\). The order identity
\eqref{cap:resolvent-order} permits this substitution in
\eqref{cap:bl-recursion}. To see the effect of the negative rank, write
\(\mathcal F_S(D)=D(I-SD)^{-1}\) and
\(B=A+Czz^\top+\eta_{\rm err}I\). Sherman--Morrison gives the exact formula
\[
 \mathcal F_S(B-d_*ee^\top)
 =\mathcal F_S(B)-
   \frac{v_Bv_B^\top}
        {d_*^{-1}+e^\top S(I-BS)^{-1}e},
 \qquad v_B=(I-BS)^{-1}e.
\]
The matrices lie in the positive domain specified above.
Since \(\|B\|=O(\delta)+\eta_{\rm err}\) and \(\|S\|=O_\gamma(1)\),
\(\mathcal F_S(B)=A+Czz^\top+O_\gamma(\delta^2)+\eta_{\rm err}I\),
\(v_B=e+O_\gamma(\delta)+\eta_{\rm err}\), and the denominator differs
from \(d_*^{-1}+e^\top Se\) by \(O_\gamma(\delta)+\eta_{\rm err}\).
After choosing the tolerances, the resulting upper bound retains the
subtraction
\begin{equation}
          -(1-\eta_{\rm small})\,\frac{ee^\top}{1/d_*+e^\top S e}
\end{equation}
up to \(O_\gamma(\delta^2)+\eta_{\rm err}\) in operator norm.

Put \(m=R_pY\) and \(T=R_c-R_p\). On \(H_p\), \(R_p\) is scalar, so
\(e=m/\|m\|\). Under the canonical bridge, \(m\) has mean \(a_c\) and
covariance \(T\). The shell likelihood and its complementary moment bounds
transfer the corresponding linear and quadratic tail tests to
\(\mathsf T\), with arbitrarily small errors in scaled units. In particular
on the shell \(\|m\|^2=(b+o(1)+O(\sigma)+O(\tau))np\),
\(\E_{\mathsf T}m\) is arbitrarily
close to \(a_c\), and
\(\E_{\mathsf T}m^\top Tm\) is at most
\(a_c^\top Ta_c+\Tr T^2\) plus an arbitrarily small scaled error.
The linear second-moment bounds are used separately for each tested
direction in \(U\), with numerical constants uniform in that direction.

At \(\alpha=0\), \(S=p^2T\). The denominator of the rank term after replacing
\(e\) by \(m/\|m\|\) is
\(\|m\|^2/d_*+p^2m^\top Tm\). The following calculation bounds its
expectation and hence the radial subtraction.

Since $R_p=p^{-2}I$ on the support of $T$,
\begin{align}
 \Tr T&=(b+o(1))n(p-c),\nonumber\\
 \Tr T^2&=(2b+o(1))n(1/c-1/p)
             -2p^{-2}(b+o(1))n(p-c).\label{cap:trace-difference-square}
\end{align}
The rank denominator before the factor $x^{-2}$ is
$\|m\|^2/d_*+p^2m^\top Tm$. Its expectation is bounded above by
\begin{equation}
 \frac{bnp}{d_*}
 +p^2\bigl(a_c^\top T a_c+\Tr T^2\bigr)+o(n p^2/c).
 \label{cap:radial-denominator}
\end{equation}
For each test vector $v$, Cauchy--Schwarz in the form $\E[(v^\top
m)^2/Q]\ge(v^\top\E m)^2/\E Q$ turns this into a rank-one matrix bound.
Dividing the denominator by $bn p^2/c=x^{-2}\|a_c\|^2(1+o(1))$ gives
successively $x/d_*$, $1-x^2$, and $2(1-x)^2$. This proves
that the radial deficit in output units is at least
\begin{equation}\label{cap:radial-regeneration}
 \frac{1-\eta_{\rm small}}
      {x/d_*+(1-x^2)+2(1-x)^2+\eta_{\rm small}}.
\end{equation}
In the stated interior-parameter regimes, \(x\ge\gamma^2\) and
\(|\alpha|\le\sigma p^2\), so
\(|\alpha|/c^2\le\sigma\gamma^{-4}\). The high-power choice of
\(\sigma\) therefore includes the resulting perturbation in
\(\eta_{\rm small}\).

Cross-direction losses in making this a purely radial deficit can be
arbitrarily small by the transverse mean and second-moment bounds above.
Formula \eqref{cap:radial-regeneration} exceeds 0.26 throughout, and exceeds
0.31 at \(x=\gamma\). The small constant losses can be chosen to keep these
strict inequalities. Remaining terms in the radial block are negligible for
this estimate.

\subsection{Compression under the fresh refinement}\label{cap:compression-proof}
For the inductive transverse block at \(x=\gamma\), first fix the complete
orientation \(O\) and a child field. Expectations in this paragraph are under
its target transition \(\mathsf T_O\). Fix a smooth function
\(0\le\chi_0\le1\), equal to one on \([-1/2,1/2]\) and zero outside
\((-1,1)\), and use
\begin{equation}\label{cap:field-cutoff}
 \chi_h(y)=\chi_0\!\left(
       \frac{\alpha_p(y)-\alpha_c(h)}{\tau p^2}\right).
\end{equation}
Here \(\tau\) is a sufficiently small fixed power of \(\gamma\), and
the support lies inside agreement with fixed slack. At the current child
field write \(\chi=\chi_h\). Write
\(A_{\rm par},z_{\rm par}\) for the parent quantities denoted by \(A,z\)
above, and set \(A=\chi A_{\rm par}\), \(z=\chi z_{\rm par}\) in the
compression estimates. The shell moment bounds make this truncation
arbitrarily small. Put
\[
 g(y)=\nabla_{H_p}f_p(y)-(p^2+\alpha)^{-1}y,\qquad
 Q_s=\mathbf1_{\{\Delta_E\ge sp^2\}} .
\]
On the support of \(\chi\),
\(\nabla_{H_p}g=D^\circ/p^2\) and
\(\|g\|\le\varepsilon_{\rm mac}M_p\), where the frozen-saddle error is
included in \(\varepsilon_{\rm mac}\). With the score
\(\mathfrak s=-\nabla_E\log\mathsf T_O\), integration by parts gives the
matrix identity
\begin{equation}\label{cap:nuclear-ibp}
 \E_{\mathsf T_O}[\chi D^\circ Q_s]
 =p^2\E_{\mathsf T_O}
    [g(\chi Q_s\mathfrak s-Q_s\nabla\chi)^\top].
\end{equation}
The nuclear norm of this expectation is bounded by
\[
 p^2\bigl(\E_{\mathsf T_O}\mathbf1_{\operatorname{supp}\chi}\|g\|^2\bigr)^{1/2}
       \bigl(\E_{\mathsf T_O}
          \|\chi Q_s\mathfrak s-Q_s\nabla\chi\|^2\bigr)^{1/2}.
\]
Summing the localized score identity used in
\eqref{cap:localized-fisher} over an orthonormal basis of \(Q_sE\) gives
\[
 \E\|\chi Q_s\mathfrak s-Q_s\nabla\chi\|^2
 =\E\chi^2\Tr(Q_sMQ_s)+\E\|Q_s\nabla\chi\|^2
 \le C_s d_p/p^2+C_s/(p^4M_p^2)
 \le C_sM_p^2.
\]
Here \(\Delta_E^{-1}\le(sp^2)^{-1}\) on \(Q_sE\), the invariant bounds
\(-K_p\) by \(C/p^2\) on the cutoff support, and
\(\|\nabla\chi\|\le C/(p^2M_p)\) with fixed buffer constants.
Thus \eqref{cap:nuclear-ibp} has nuclear norm at most
\(C_s\varepsilon_{\rm mac}p^2M_p^2
=C_s\varepsilon_{\rm mac}k_p\).

The columns omitted by \(Q_s\) have rank \(O(s)k_p+o(k_p)\), by the edge
count in the strip \(0\le\Delta<sp^2\), including the zero-increment
boundary coordinates; their operator norm is bounded.
Changing \(D^\circ\) to \(D\) costs \(C\tau d_p\) in nuclear norm on the
shell. Finally, \(D-A_{\rm par}\) has rank at most two for each field and bounded
nuclear norm. Its radial direction may change with the field, but
\[
 \|\E_{\mathsf T_O}[\chi(D-A_{\rm par})]\|_*
 \le\E_{\mathsf T_O}\|\chi(D-A_{\rm par})\|_*\le C.
\]
This uses the nuclear triangle inequality, without assigning a fixed rank
to the average. Choosing \(s\) first, then the macroscopic and shell errors,
and finally \(k_p\) large proves
\begin{equation}\label{cap:nuclear-average}
 \frac{\|\bar A_O\|_*}{k_p}\le\eta_{\rm err},
 \qquad \bar A_O=\E_{\mathsf T_O}A .
\end{equation}
The norm is outside the expectation. In particular the signed normalized
trace is small.

The radial correction has the required sign, but a direct operator-norm
estimate on the transverse block would be enlarged by the change of scale. The
small signed average just proved allows a better estimate. We sample finitely
many independent terminal fields, show that their average has small Frobenius
norm, and compress it using the remaining Haar rotation. The sampled field
directions are exceptional because each matrix depends on its own field; we
estimate their span separately.

We need
\begin{equation}
 \|[\E_{\mathsf T_O}A]_{U_\perp U_\perp}\|
 \le o_{\gamma\to0}(1)\,\gamma^2\delta,\qquad
 \E_{\mathsf T_O}\|P_U z\|^2\le C K_0^2\delta\, H\gamma^3 ,
 \label{cap:compression}
\end{equation}
apart from tiny terms, where \(H=C_1\log(1/\gamma)\) with a sufficiently large
fixed coefficient. These assertions hold simultaneously for child fields with
high probability over the refinement. For the moment, keep \(O\) and the child
field fixed and sample \(Y_1,\ldots,Y_l\) independently from \(\mathsf T_O\).
The deterministic function of these fields is
\(\mathcal A(Y_1,\ldots,Y_l)=l^{-1}\sum_jA(Y_j)\). Independence gives exactly
\begin{align}\label{cap:sample-frobenius}
 \E_{\mathsf T_O^l}\|\mathcal A\|_F^2
 &=l^{-1}\E_{\mathsf T_O}\|A\|_F^2
   +\frac{l-1}{l}\|\bar A_O\|_F^2\nonumber\\
 &\le \delta^2d_p/l+\delta\|\bar A_O\|_* .
\end{align}
The last inequality uses \(\|A\|\le\delta\) and
\(\|\bar A_O\|_F^2\le\|\bar A_O\|\|\bar A_O\|_*\).
By \eqref{cap:nuclear-average}, the quotient of
\eqref{cap:sample-frobenius} by \(k_p\) is
\(O(\delta^2/l)+\eta_{\rm err}\). We will combine this fixed-\(O\) target-law
estimate with a Haar event first proved under canonical sampling.

\paragraph{Conditional Haar laws before tilting.}
We spell out the conditional laws because the actual transition is not itself
Haar after its terminal fields have been sampled. Fix $\mathscr F_p$ and a
child field in spectral coordinates, denoted $u\in U_0$. Its image is
$h_c=Ou$, and $U=OU_0$. Sample $l$ independent canonical terminal fields
$\xi_1,\ldots,\xi_l$ in spectral coordinates according to
\eqref{cap:canonical-bridge}, independently of $O$. Their means are parallel
to \(u\). Denote this joint law by
\[
 \mathbb P_{\rm can}(dO,d\xi_1\cdots d\xi_l)
 =\operatorname{Haar}(dO)\prod_{j=1}^l
       \mathsf G_{u,\rm diag}(d\xi_j).
\]
Its physical fields are \(Y_j=O\xi_j\), whose conditional law for fixed
\(O\) is \(\mathsf G_O^l\). Conditional on these spectral fields, and on the images under \(O\)
of $u$ and any selected collection of the $\xi_j$, the remaining map on the
orthogonal complements is Haar. Conditional on the image of $\xi_j$, the
matrix $A_j=A(O\xi_j)$ is fixed: the parent potential, its radial splitting,
and its derivatives are all measurable with respect to $\mathscr F_p$ and
that image. Conditional on two images, $A_iA_j$ is fixed. The same statement
applies to $z_j=z(O\xi_j)$, with $z_j\perp O\xi_j$.

Set $e_j=O\xi_j/\|\xi_j\|$, and let $e_0=h_c/\|h_c\|$. The child annulus is
separated from zero field, so $e_0$ is defined. All dimensions below are
comparable to $k_p$ or $k_c$ by the spectral count. The elementary Haar
estimate we use is
\begin{equation}\label{cap:haar-compression}
 \|P_V O^\top T O P_V-\tfrac{\Tr T}{d}P_V\|
 \le C\left(\frac{\|T\|_F}{d}\sqrt{r+a}
                 +\frac{\|T\|}{d}(r+a)\right),
\end{equation}
with failure probability at most $2e^{-a}$, when $O$ is Haar in dimension $d$,
$V$ is a fixed $r$-dimensional subspace, and $T$ is symmetric and fixed.
Replacing $r+a$ by a fixed larger multiple absorbs the size of a
fixed-resolution unit net in $V$. To prove the estimate for a fixed unit
vector, write its Haar image as $g/\|g\|$, diagonalize $T$, and use
\begin{equation}
 \log\E e^{t\sum_i\lambda_i(g_i^2-1)}
 \le C t^2\sum_i\lambda_i^2,
 \qquad |t|\max_i|\lambda_i|\le\tfrac14.
 \label{cap:gaussian-quadratic-tail}
\end{equation}
The same bound for $\|g\|^2$ handles its normalization. Polarization gives
the bilinear and nonsymmetric versions used for products.

On the complement of the span of $e_0,e_1,\ldots,e_l$, apply
\eqref{cap:haar-compression} with $T=\mathcal A$ and $a$ a sufficiently large
multiple of $Hk_c$. The deviation from its scalar trace is at most
\begin{equation}\label{cap:ordinary-compression}
 C\left(\frac{\|\mathcal A\|_F}{\sqrt{k_p}}
                      \sqrt{H\gamma^3}
             +\delta H\gamma^3\right).
\end{equation}
Conditioning only on the image of $\xi_j$, the vector $z_j$ is fixed and
perpendicular to $e_j$. The residual projection of $U$ has rank at most
$\dim U$, so the same Haar estimate gives
$\|P_Uz_j\|^2\le C K_0^2\delta H\gamma^3$.
Its failure rate can be $2Hk_c$ after increasing absolute constants.

\paragraph{The span of sampled directions.}
Write $W$ for the span of the $e_j$ after projection onto $e_0^\perp$.
Under canonical sampling, the centered Gaussian vectors before direction
normalization have squared norm comparable to one in field units, limiting
orthogonal Gram matrix, $U_0$-projection squared norm of order $\gamma$,
and covariance operator norm at most $C/(\gamma^2k_p)$. Gaussian quadratic
tails and nets in the fixed $l$-dimensional coefficient space therefore give
\begin{equation}\label{cap:exceptional-span-geometry}
 \|P_{U\cap e_0^\perp}P_W\|\le C\sqrt{H\gamma},\qquad
 \sum_{j=1}^l |\langle e_j,w\rangle|^2\le C\|w\|^2
 \quad(w\in W).
\end{equation}
The projected generating vectors have a synthesis map and an inverse on its
range bounded by absolute constants. Any pair of the original directions has
bounded Gram conditioning. These conclusions hold at the same $2Hk_c$ rate.
Although $l$ will be a large power of $\gamma^{-1}$, it is fixed before
$n\to\infty$, so the fixed-dimensional nets cost no additional multiple of
$k_p$ in the limit.

Let \(\bar a=\Tr\mathcal A/d_p\). Then
\begin{align}
 \|P_W(\mathcal A-\bar a I)P_W\|
 &\le C\delta\bigl(\sqrt{H\gamma^3}+l^{-1/2}\bigr),
 \label{cap:exceptional-first-moment}\\
 \|P_W\mathcal A^2P_W\|
 &\le C\|\mathcal A\|_F^2/k_p
       +C\delta^2\bigl(\sqrt{H\gamma^3}+l^{-1/2}\bigr).
 \label{cap:exceptional-second-moment}
\end{align}
To account for the dependence of each matrix on its own sampled direction,
expand a test vector in the bounded generating frame of \(W\). For
the summand $A_j$, condition on $e_0,e_j$ and apply
\eqref{cap:haar-compression} to the other generating directions in their
orthogonal complement. The scalar term is $\Tr A_j/d_p$ up to a
fixed-dimensional trace correction, which is $o(1)$. Terms with a row or
column in $\operatorname{span}(e_0,e_j)$ have average bounded by
\[
 \frac{C\delta}{l}\sum_j\|P_{\operatorname{span}(e_0,e_j)}w\|
 \le C\delta l^{-1/2}\|w\|,
 \qquad w\in W,
\]
using $w\perp e_0$ and \eqref{cap:exceptional-span-geometry}. This proves
\eqref{cap:exceptional-first-moment} by polarization and a fixed-dimensional
net. For the second estimate expand $\mathcal A^2=l^{-2}\sum_{i,j}A_iA_j$.
Condition on $e_0,e_i,e_j$. The same Haar test gives scalar term
$\Tr(A_iA_j)/d_p$; the sum of these terms is $\Tr\mathcal A^2/d_p$. The
exceptional-row contribution is bounded by $C\delta^2l^{-1/2}$, since
\[
 l^{-2}\sum_{i,j}
  \|P_{\operatorname{span}(e_0,e_i,e_j)}w\|^2
 \le C l^{-1}\|w\|^2.
\]
The same bound includes the diagonal pairs. Summing the scalar terms and the
exceptional-row errors proves \eqref{cap:exceptional-second-moment}. Before
applying the quadratic Haar test to a product, symmetrize it or use
polarization. The argument therefore does not require an individual $A_iA_j$
to be positive.

Conditional on all the sampled images, the residual rotation gives a factor
$C\sqrt{H\gamma^3}$ on the nonexceptional side of the block from $W$ to its
complement. Together with \eqref{cap:exceptional-span-geometry}, this is the
product $C\sqrt{H\gamma}\sqrt{H\gamma^3}$, not merely the first factor. There
is no child radial contribution, since the target tests lie in $U\cap
e_0^\perp$. Combining the three blocks gives, in addition to the scalar trace,
\begin{align}
 C\biggl[&
 \frac{\|\mathcal A\|_F}{\sqrt{k_p}}\sqrt{H\gamma^3}
 +\delta H\gamma^3
 +\delta H\gamma\bigl(\sqrt{H\gamma^3}+l^{-1/2}\bigr)\nonumber\\
 &+H\gamma^2
 \left(\frac{\|\mathcal A\|_F}{\sqrt{k_p}}
       +\delta(H\gamma^3)^{1/4}+\delta l^{-1/4}\right)
 \biggr].\label{cap:compression-remainder}
\end{align}
Choose \(l=\lceil\gamma^{-16}\rceil\). Up to this point,
\eqref{cap:compression-remainder} is a deterministic bound on the function
\(\mathcal A(Y_1,\ldots,Y_l)\) outside a bad event under the joint canonical
law \(\mathbb P_{\rm can}\). The fixed-\(O\) estimates
\eqref{cap:nuclear-average} and \eqref{cap:sample-frobenius} concern
\(\mathsf T_O^l\). We transfer the geometric event to that law before using
those expectations.

\paragraph{Transfer to actual transitions and a field net.}
So far the Haar estimates concern canonical Gaussian samples. We next
convert them into estimates for the actual transition while keeping track
of the loss in exponential rate. On a narrow parent shell the density ratio
costs an arbitrarily small fixed multiple of the parent dimension for each
sampled field. We choose that loss only after fixing the sample count, so
the remaining rate pays for the child field net. Shell complements already
have uniform conditional bounds and are added separately.


For fixed \(O\), let \(F_O\) denote \(V_p-q_{p,\alpha}^{\rm fr}\) in
\eqref{cap:bridge-density-ratio}. Its normalizing denominator may depend on
the remaining orientation even after the terminal images have been fixed.
Thus the residual Haar statement above is used only under
\(\mathbb P_{\rm can}\). On the shell,
\eqref{cap:shell-likelihood} bounds each likelihood ratio by \(e^{\eta k_p}\).
For the canonical bad event \(\mathcal E\), the joint change of law gives
\[
 \E_O\mathsf T_O^l(\mathcal E\cap\mathcal S_\tau^l)
 \le e^{l\eta k_p}\Pr_{\rm can}(\mathcal E).
\]
Choose \(\eta\ll H\gamma^3/l\). A canonical bound \(e^{-2Hk_c}\)
then makes the right-hand side at most \(e^{-3Hk_c/2}\).
Markov's inequality in \(O\) gives the explicit conditional form
\[
 \Pr_O\left\{
 \mathsf T_O^l(\mathcal E\cap\mathcal S_\tau^l)>e^{-Hk_c/2}
 \right\}\le e^{-Hk_c}.
\]

The complement of the narrow shell is handled separately. Its conditional
probability and all bounded polynomial moment errors are exponentially small,
uniformly at every child field, by the deterministic upper support bound and
Gaussian distance tails already proved. This uniform estimate is not obtained
by a union bound over orientations or field points, and its exponent need not
be \(2Hk_c\). The scalar signed-trace and sample independence calculations
are also applied anew conditionally at each field.

For an orientation outside the displayed exceptional set, we may now take
the \(\mathsf T_O^l\)-expectation of \eqref{cap:compression-remainder}.
Equation~\eqref{cap:sample-frobenius} bounds the expected Frobenius term by
\(O(\delta/\sqrt l)+\eta_{\rm err}\), and
\eqref{cap:nuclear-average} bounds the expected scalar trace.
The bad-event and shell-complement contributions are small because the
matrices have bounded operator norm there. After division by
\(\gamma^2\delta\), the displayed remainder terms are bounded by constant
multiples of \(H\gamma\), \(H^{3/2}\gamma^{1/2}\),
\(H^{5/4}\gamma^{3/4}\), and powers of \(\gamma\) from
\(l^{-1/4}\) and \(l^{-1/2}\). They tend to zero as \(\gamma\to0\).
This proves the first estimate in \eqref{cap:compression}, with an
arbitrarily small additive error. Applying the same change of law to the
one-vector event gives the stated bound for
\(\E_{\mathsf T_O}\|P_Uz\|^2\).

We extend this estimate to all fields by comparing geometric bad events,
before taking their signed matrix expectations. Take a net of resolution
\(\rho\) in child field units, with \(\rho\) a sufficiently high fixed
power of \(\gamma\). Its size is at most
\(\exp\{C k_c\log(1/\gamma)\}\). Fix an orientation and neighboring
fields with \(\|h-h'\|\le\rho c^2M_c\), and keep the physical terminal
samples fixed. Take \(\rho\) small enough that their joining segment
still has its frozen parameter inside parent agreement. The untruncated
parent matrices are then unchanged. The
scaled saddle Lipschitz bound gives
\begin{equation}\label{cap:cutoff-continuity}
 \sup_y|\chi_h(y)-\chi_{h'}(y)|
       \le C\gamma^2\rho/\tau.
\end{equation}
Thus the cutoff-weighted \(A\) changes by at most
\(C\delta\gamma^2\rho/\tau\) in operator norm, and its Frobenius norm
divided by \(\sqrt{k_p}\) changes by the same order. The corresponding
bound for \(z\) has \(K_0\sqrt\delta\) in place of \(\delta\).
The child radial projection changes by \(C\rho\), since the annulus is
separated from zero. On the retained Gram event the generating frames and
their inverses have bounded norms. Consequently the geometric inequalities
leading to \eqref{cap:compression-remainder}, with fixed slack in their
constants, persist under these changes with an additional
\(C_\gamma\rho\delta\) error. The same holds for the one-vector test.
The constant \(C_\gamma\) is at most a fixed power of \(\gamma^{-1}\):
the sample count, cutoff width, and frame bounds have already been fixed.

The transition densities require an exponential comparison. For a fixed
orientation, on \(\|Y\|\le C p^2M_p\),
\begin{equation}\label{cap:field-likelihood-comparison}
 \left|\log\frac{d\mathsf T_h}{d\mathsf T_{h'}}(Y)\right|
       \le C_\gamma\rho k_c.
\end{equation}
Indeed the potential \(V_p(Y)\) cancels in the ratio, and the Gaussian
log ratio uses only \(\Delta_U^{-1}\). The normalizing factor obeys the
same bound: differentiating the convolution gives
\(\nabla_U\phi_c=\Delta_U^{-1}(\E_{\mathsf T_h}Y_U-h)\),
whose norm is \(O_\gamma(M_p)\) by the parent shell and complementary
moment bounds. Integrating this derivative along the short field segment
proves the normalization assertion. This argument uses only the conditioned
parent hypotheses, not a separate good event at the child.

Choose \(\rho\) so that the perturbation of the matrix tests is much
smaller than \(\gamma^2\delta\), and
\(lC_\gamma\rho<H/8\). Prove the net estimates with the stated strict
slack. Outside their orientation exceptions, the conditional probability of
a relaxed geometric failure at a neighboring field, with all samples in a
narrower typical shell, is at most
\(e^{lC_\gamma\rho k_c}e^{-Hk_c/2}\le e^{-3Hk_c/8}\).
The cutoff and radial continuity above justify this event inclusion; the
narrower shell stays in the shell used at the net field. The shell
complements and large lengths have the separate uniform conditional moment
bounds already proved. At the neighboring field apply
\eqref{cap:nuclear-average} and \eqref{cap:sample-frobenius} anew, and then
take the expectation of the relaxed geometric inequality. This retains
their signed cancellation and proves \eqref{cap:compression} everywhere
on the annulus.

Choose the fixed exponents in \(\tau,\rho\) and the sample count first,
then \(C_1\) in \(H=C_1\log(1/\gamma)\) to pay for the net, and then
\(\gamma\) small. The separate uniform shell bounds are added without a
field union bound. The resulting failure probabilities have the form
\(e^{-a k_p}+e^{-a'k_c}\) with fixed positive constants and are summable
on the geometric list down to \(k_p\asymp\sqrt{\log n}\).

\subsection{Completion of one inductive step}
We can now assemble the matrix estimates. Compression makes the
transverse contributions in both transport bounds smaller than the
allowed transverse error. Radial regeneration supplies strict slack in the
upper radial block, and the Fisher test supplies a bounded lower radial
block. Positive-matrix Cauchy--Schwarz between these two bounds then controls
the actual mixed block. This closes the four parts of the invariant at the
child scale. The intermediate-scale estimates use the same deterministic
transport calculations; they do not require another random compression.

By \eqref{cap:compression}, \eqref{cap:bl-recursion} has transverse surplus
\(o(1)\delta\), since the mixed-square term gains \(\gamma^3 H\). For the
Fisher lower bound, let \(e_0\) be the child radial direction. The canonical
mean \(a_c\) is parallel to \(e_0\), and
\(v^\top(R_c-R_p)v\le C_\gamma p^{-2}\) for each unit
\(v\in U\cap e_0^\perp\). The shell norm and transferred second moments
from the radial calculation therefore give
\[
 \sup_{\substack{v\in U\cap e_0^\perp\\\|v\|=1}}
       \E_{\mathsf T}(v^\top e)^2
       \le \eta_{\rm err}+C_\gamma/k_p.
\]
The bounded radial entry of \(D\) is negligible on these tests. Its mixed
terms are negligible by Cauchy--Schwarz and \(\|z\|\le K_0\sqrt\delta\),
after choosing the comparison error and then \(k_p\). The transverse
compression thus gives a Fisher deficit \(o(1)\delta\) as well, including
the fixed factor \(x^2N^2\). The transverse-only Fisher test gives a block lower matrix
with bounded radial deficit, and the Brascamp--Lieb upper matrix has cross
error \(O_\gamma(\delta)+\eta_{\rm err}\) after taking a purely radial gain as
above. The two Loewner bounds therefore force an actual cross block of size
\(O(\sqrt{o(1)\delta})+O_\gamma(\delta)\), by positive-matrix Cauchy--Schwarz.
This proves \eqref{cap:invariant} inductively, with slack. At intermediate
scales in the stated interior-parameter regime the same calculations without
\eqref{cap:compression} give transverse error and actual mixed block bounded
by a quantity tending to zero (at fixed \(\gamma\)), bounded total Hessian,
and radial deficit at least 0.26. For the lower mixed estimate near \(x=1\)
use the full Fisher test; its predicted cross block is small by the same
bridge projections. These intermediate bounds are deterministic and uniform
under the parent hypotheses.

\subsection{Deterministic evaluation at other cap ratios}

Once a parent satisfies the invariant, evaluation at another scale uses
its deterministic support and covariance bounds. No transverse compression
is needed if the target is an upper curvature estimate with prescribed
fixed error, rather than a renewed invariant with the same error.

\begin{lemma}[Evaluation with the induction extension]
\label{cap:deterministic-evaluation}
Under the parent hypotheses of Proposition~\ref{cap:curvature-induction},
the deterministic intermediate-scale conclusions hold also for
\(x=c/p\in[\gamma^2,\gamma]\), whenever the frozen parameter is inside
parent agreement with fixed slack. At fixed \(\gamma\), the transverse
and mixed errors tend to zero with the parent errors, the scaled Hessian
is bounded, and the radial deficit is at least \(0.26\).
\end{lemma}
\begin{proof}
The same extension has a positive convolution margin because
\[
 \|\Delta^{1/2}L_p\Delta^{1/2}\|
 \le(1+\gamma^4/100)(1-\gamma^4)<1.
\]
All inverse bounds in the shell estimates and the upper and lower
covariance transports remain bounded by fixed powers of
\(\gamma^{-1}\). Thus those deterministic estimates apply with enlarged
constants. Formula~\eqref{cap:radial-regeneration} still gives the stated
radial deficit, while the other blocks tend to zero with the input errors.
The proof uses no additional orientation event and does not extend the
probabilistic induction statement beyond ratio \(\gamma\).
\end{proof}

For a bounded child field ball, the saddle parameter can occupy a wider
range in child units. We choose a coarser parent and center a new support
annulus at that frozen parameter. This is a separate evaluation of the
true potential; its extension is not propagated as the next inductive
potential.

\begin{lemma}[Supports centered at a frozen parameter]
\label{cap:frozen-evaluation-support}
Fix a child field radius \(R\) before choosing the curvature and
comparison tolerances. Suppose a true parent at scale \(p\) satisfies
the invariant and the static comparisons. There is a fixed \(b_1>0\),
depending on \(R,\sigma\), such that the following holds for every
\(x\in[\gamma b_1,b_1]\), \(c=xp\), and
\(h\in H_c\) with \(\|h\|\le R c^2M_c\). Put
\(\alpha=\alpha_c(h)\). With sufficiently small fixed tolerances, there
is an upper-semiconcave extension agreeing with the true parent on a
neighborhood of the frozen typical length sphere. Its convolution has a
strict margin, and its terminal mean and covariance differ from those of
the true transition by exponentially small errors in the corresponding
scaled units. Consequently,
\begin{equation}\label{cap:ball-evaluation-bound}
 K_c(h)\preceq C(R)c^{-2}I_{H_c},
\end{equation}
where \(C(R)\) is independent of the chosen ratios and sharp curvature
tolerances. The construction and estimates are uniform on the indicated
fixed ratio interval, with \(n\) taken large after those choices.
\end{lemma}
\begin{proof}
The static saddle bounds give \(\kappa_R,C_R>0\), depending only on
the field radius, such that
\[
 -1+\kappa_R\le\alpha/c^2\le C_R.
\]
Choose \(b_1\) so that
\(\max(1,C_R)b_1^2<\sigma/2\). Then \(|\alpha|<\sigma p^2/2\)
throughout the ratio interval. Choose fixed numbers \(e,w>0\), also
satisfying the tail and scalar-margin bounds below, with
\[
 C(w+\delta)\ll e\ll\kappa_R(\gamma b_1)^2,
 \qquad w<\sigma/4,
\]
and define nested shells
\[
 S_j^\alpha=\{y\in H_p:|\alpha_p(y)-\alpha|\le t_jwp^2\},
 \qquad (t_0,t_1,t_2)=(1/2,3/4,1).
\]
They lie inside the parent invariant annulus. The radius derivative
computed above gives buffers comparable to \(wp^2M_p\) between them.
Set
\[
 L_\alpha=\bigl((p^2+\alpha)^{-1}+ep^{-2}\bigr)I_{H_p}.
\]
On \(S_2^\alpha\), the invariant's positive-part bound gives
\[
 K_p(y)\preceq
 \bigl((p^2+\alpha)^{-1}+C(w+\delta)p^{-2}\bigr)I_{H_p}
 \prec L_\alpha.
\]
This proves the upper-support inequality for pairs in \(S_1^\alpha\)
whose joining segment remains in \(S_2^\alpha\).

For longer pairs, freeze the own saddle at the support point and use
\eqref{cap:frozen-mgf}, the value comparison, and the gradient comparison
as in \eqref{cap:parent-upper-support}. If the buffer divided by
\(p^2M_p\) is \(\beta\asymp w\), the excess over the frozen
\(\alpha\) quadratic is at most
\[
 \frac{\|y-z\|^2}{2p^2}
 \left[Cw+
 2\left(\frac{2\varepsilon_v+C_{\rm dens}/k_p}{\beta^2}
               +\frac{\varepsilon_g}{\beta}\right)\right].
\]
Take \(\varepsilon_v+C_{\rm dens}/k_p\ll ew^2\) and
\(\varepsilon_g\ll ew\). The bracket is smaller than \(e\), so the
upper-support inequality holds for every pair in \(S_1^\alpha\).
Define
\[
 \bar\phi_{p,\alpha}(y)=\inf_{z\in S_1^\alpha}
 \left\{\phi_p(z)+\nabla_{H_p}\phi_p(z)^\top(y-z)
                   +\tfrac12(y-z)^\top L_\alpha(y-z)\right\}.
\]
The argument of \eqref{cap:parent-extension} shows that this function
agrees with \(\phi_p\) on \(S_1^\alpha\), has at most quadratic
growth, and satisfies \(\nabla^2\bar\phi_{p,\alpha}\preceq L_\alpha\)
distributionally. Its dependence on the child field is through the scalar
\(\alpha\), determined by the eigenvalues and \(\|h\|\). No child
angular refinement enters its construction. For this evaluation hold the
selected extension fixed when differentiating its convolution at \(h\).
Different fields may select different extensions; no derivative of this
selection, or of its parameter \(\alpha\), enters the Hessian formula.

We verify the complementary moments before comparing covariances. Let
\(B_\alpha\) denote the logarithm of the frozen norm density at a field
on the sphere \(\|y\|=\ell_p(\alpha)\). It is constant on that sphere.
For \(d(y)=|\|y\|-\ell_p(\alpha)|/(p^2M_p)\), the nearest support
point on the sphere gives
\[
 \bar\phi_{p,\alpha}(y)-q_{p,\alpha}^{\rm fr}(y)
 \le B_\alpha+k_p\bigl(\varepsilon_v+
                    \varepsilon_gd(y)+\tfrac12ed(y)^2\bigr)
 \le B_\alpha+k_p\bigl(\varepsilon_v+
                    \varepsilon_g^2/e+ed(y)^2\bigr).
\]
The canonical calculation leading to \eqref{cap:bridge-distance-tail}
now gives, with constants depending only on the bounded child saddle range,
\[
 \Pr_{\mathsf G}\{d(Y)>v\}\le2e^{-c_Rx^2k_pv^2},
 \qquad v\ge C_R/(x\sqrt{k_p}).
\]
Require also \(e<c_R(\gamma b_1)^2/4\). Normalize the
extended transition on a shell
\(|\alpha_p(Y)-\alpha|\le\tau_0p^2\) with
\(\tau_0\ll\gamma b_1w\). Its canonical probability tends to one.
The off-saddle identity \eqref{cap:off-saddle-density} and centered density
bounds give a logarithmic normalization loss at most
\(\varepsilon_vk_p+C\tau_0^2k_p+O(1)\) relative to \(B_\alpha\).
Choose the comparison errors and then \(n\) so that
\[
 2\varepsilon_v+\varepsilon_g^2/e+C\tau_0^2+O(1)/k_p
       \ll (\gamma b_1)^2w^2.
\]
Leaving \(S_1^\alpha\) requires \(d(Y)\ge c w\). Integrating the
Gaussian distance tail against the upper support therefore bounds every
fixed polynomial moment of \(1+d(Y)\) on this complement by
\(C_j e^{-a k_p}\), for a fixed \(a>0\), uniformly in \(x,h\).

The true transition needs its own complementary estimate. On each fixed
bounded parent field range, static comparison and the frozen numerator
density upper bound give
\[
 \phi_p(y)-q_{p,\alpha}^{\rm fr}(y)
       \le B_\alpha+\varepsilon_vk_p+C_{\rm dens}.
\]
The preceding normalization shell and Gaussian distance tail control the
bounded part of the complement. For the remaining tail, sample the child
spin by tilting the parent zero-field spin law with
\(\exp\{h^\top X+X^\top\Delta X/2\}\), then set
\(Y=h+\Delta X+N(0,\Delta)\). The tilt normalizer is at least one by
zero-field symmetry and \(\Delta\succeq0\). With
\(v=\|P_pX\|/M_p\), its numerator is at most
\(e^{C_R k_p(1+v^2)}\). The parent projection tail
\(e^{-ck_pv^6}\) survives this factor, and the observation noise has
a Gaussian tail in parent field units. Increasing the fixed outer radius
therefore makes the outer contribution and all its fixed polynomial
moments exponentially small. This argument uses the true parent projection
tails, not the extension's Hessian bound.

The two unnormalized transition laws agree on \(S_1^\alpha\). Their
normalization ratio and their complementary first and second moments thus
show that their means and covariances differ exponentially in parent scaled
units, as in the earlier transition comparison. Converting a covariance
error in units \((p^2M_p)^2\) to child Hessian units multiplies it by
\[
 c^2\|\Delta_U^{-1}\|^2(p^2M_p)^2
       =\frac{x^2k_p}{(1-x^2)^2}.
\]
The exact formula \eqref{cap:child-hessian} therefore leaves an error in
\(c^2K_c\) tending to zero uniformly on the fixed ratio interval.

Finally put \(a=\alpha/c^2\) and
\(q=(1+x^2a)^{-1}+e\). On \(U\), the increment is
\((1-x^2)p^2I\), while the global support curvature is \(q/p^2\).
Impose also
\(e\le\kappa_R(\gamma b_1)^2/[4(1+\sigma/2)]\).
Then the convolution has a strict margin. Brascamp--Lieb, first for smooth
approximations and then by passage to the limit, and
\eqref{cap:child-hessian} give
\[
 c^2K_c^{\rm ext}
 \preceq
 \frac{1+(1+x^2a)e}
      {(1+a)-(1-x^2)(1+x^2a)e/x^2}\,I_U
 \preceq \frac{8}{3\kappa_R}I_U.
\]
The denominator is at least \(3\kappa_R/4\), and the numerator is at
most two. The exponentially small replacement error is absorbed by
enlarging this constant once. This proves \eqref{cap:ball-evaluation-bound}
with no dependence of its coarse constant on \(x\) or the sharp
tolerances. All shell widths, comparison errors, and the ratio interval are
fixed before \(n\to\infty\); no estimate here requires taking
\(x\to0\) while keeping these choices fixed.
\end{proof}
